\begin{figure}[H]
\centering
\begin{tikzpicture}[>=Stealth,font=\small]
  \def\sc{2.0}   % cm por unidad de  xi*hbar^2

  % eje de energía
  \draw[->,grisUNI] (-0.8,-2.9)--(-0.8,2.0) node[above]{$E$};

  % nivel sin acoplamiento
  \draw[very thick] (0,0)--(2.2,0);
  \node[anchor=east] at (-0.05,0) {$2p$};
  \node[above,font=\footnotesize] at (1.1,0.05) {$g=6$};

  % niveles desdoblados
  \draw[very thick,rojoUNI] (5,{0.5*\sc})--(7.2,{0.5*\sc});
  \draw[very thick,azulUNI] (5,{-1.0*\sc})--(7.2,{-1.0*\sc});
  \draw[dotted,grisUNI] (2.2,0)--(5,{0.5*\sc});
  \draw[dotted,grisUNI] (2.2,0)--(5,{-1.0*\sc});

  \node[rojoUNI,anchor=south west,font=\small] at (5,{0.5*\sc+0.05})
       {${}^{2}\!P_{3/2}$\; ($g=4$)};
  \node[rojoUNI,anchor=north west,font=\footnotesize] at (5,{0.5*\sc-0.05})
       {$\Delta E=+\tfrac12\,\xi\hbar^2$};
  \node[azulUNI,anchor=south west,font=\small] at (5,{-1.0*\sc+0.05})
       {${}^{2}\!P_{1/2}$\; ($g=2$)};
  \node[azulUNI,anchor=north west,font=\footnotesize] at (5,{-1.0*\sc-0.05})
       {$\Delta E=-\,\xi\hbar^2$};

  % centro de gravedad
  \draw[dashed,naranjaUNI] (2.2,0)--(7.5,0);
  \node[naranjaUNI,anchor=south west,font=\footnotesize,align=left] at (7.7,0.1)
       {centro de gravedad:\\ $4\!\cdot\!\tfrac12+2\!\cdot\!(-1)=0$};

  % desdoblamiento
  \draw[<->,thick] (7.5,{-1.0*\sc})--(7.5,{0.5*\sc});
  \node[anchor=west,font=\footnotesize,align=left] at (7.7,{-0.75*\sc})
       {$\Delta E_{\rm fina}=\tfrac32\,\xi\hbar^2$};
\end{tikzpicture}
\caption{Estructura fina del nivel $2p$ ($\ell=1$, $s=\tfrac12$) por
$\hat H_{SO}=\xi\,\hat{\vec L}\cdot\hat{\vec S}$ con $\xi>0$. La perturbación
reparte el nivel en $^2P_{3/2}$ (subida) y $^2P_{1/2}$ (bajada) de modo que
$\sum_j(2j+1)\Delta E_j=0$: la traza de $\hat H_{SO}$ en el subespacio de
dimensión 6 es nula y la energía media no cambia.}
\label{fig:estructura-fina}
\end{figure}
